\section{Experiments}
\label{sec:experiments}
Our experiments examine whether execution feedback improves task completion, how data augmentation and planning correction interact, and which mechanisms account for the gains.

\subsection{Experimental setup}
\paragraph{Tasks and metrics.}
PushObj comprises 350 image-goal tasks across seven object shapes; I, small-T, and square are unseen during training. Maze comprises 100 navigation tasks across five held-out layouts. PushT evaluates 50 matched start--goal tasks under clean observations and six visual shifts: blur, salt-and-pepper noise, darker lighting, and changes to the agent, object, or goal color. Manipulation success requires joint pusher/object position error below 20 pixels and object-angle error below $\pi/9$; Maze uses a 0.5-unit position threshold.

\paragraph{Methods and planning budget.}
We compare with DINO-WM~\cite{zhou2025dinowm}, LeWM~\cite{maes2026lewm}, DDP-WM~\cite{yin2026ddpwm}, and the released world model with frozen parameters or AdaJEPA updates~\cite{wang2026adajepa}. The main comparison reports PushObj and PushT together, using a 150-control execution budget and each method's gradient descent (GD) or cross-entropy (CEM) planner. Our evaluations allow 30 MPC replans, executing five controls per plan with a 25-control horizon; Maze uses action repeat four. The released and augmented-data checkpoints are evaluated separately. Within each checkpoint, ECP-on/off comparisons share tasks and action initialization. ECP uses $S{+}P$ for manipulation and gated $S$ for Maze. Planner schedules, checkpoint provenance, and task construction are specified in \cref{sec:implementation,sec:pusht-extension}.

\subsection{Comparison with recent world models}
\begin{table*}[t]
\centering
\caption{\textbf{World-model control on PushObj and PushT.} Success (\%); higher is better. PushObj has two distinct 350-task lists: archived baseline tasks (Reference) and intervention-study tasks (Paired). Results are compared within each list. PushT uses 50 shared tasks per condition; bold marks its column best. All controllers have a 150-control budget. Agent, Object, and Goal denote red-color shifts; Avg. averages the six shifts. Ours combines data augmentation and ECP; dashes denote unavailable measurements. Protocols appear in \cref{sec:pushobj-comparison,sec:pusht-extension}.}
\label{tab:main}
\small
\setlength{\tabcolsep}{3pt}
\begin{tabular*}{\textwidth}{@{\extracolsep{\fill}}llrrrrrrrrrr@{}}
\toprule
& &\multicolumn{2}{c}{PushObj}&\multicolumn{8}{c}{PushT}\\
\cmidrule(lr){3-4}\cmidrule(lr){5-12}
Method & Opt. & Reference & Paired & Clean & Blur & S\&P & Dark & Agent & Object & Goal & Avg.\\
\midrule
DINO-WM~\cite{zhou2025dinowm} & GD & 49.1 & --- & 50 & 44 & 8 & 50 & 44 & 36 & 46 & 38.0\\
LeWM~\cite{maes2026lewm} & GD & 60.6 & --- & 64 & 62 & 62 & 58 & 4 & 16 & 34 & 39.3\\
DINO-WM~\cite{zhou2025dinowm} & CEM & 62.0 & --- & 48 & 40 & 2 & 44 & 40 & 24 & 48 & 33.0\\
LeWM~\cite{maes2026lewm} & CEM & 64.0 & --- & --- & 66 & 64 & 64 & 2 & 20 & 26 & 40.3\\
DDP-WM~\cite{yin2026ddpwm} & CEM & 57.7 & --- & --- & 56 & 22 & 66 & 62 & \textbf{42} & \textbf{68} & 52.7\\
\midrule
Frozen (released) & GD & --- & 43.7 & 60 & 50 & 60 & 52 & 56 & 10 & 36 & 44.0\\
AdaJEPA (released) & GD & --- & 69.4 & 82 & 76 & 76 & 76 & 60 & 10 & 40 & 56.3\\
\midrule
\textbf{Ours (frozen)} & GD & --- & 72.3 & 82 & 72 & 70 & 78 & \textbf{68} & 12 & 34 & 55.7\\
\textbf{Ours (AdaJEPA)} & GD & --- & 76.9 & \textbf{86} & \textbf{86} & \textbf{80} & \textbf{88} & \textbf{68} & 18 & 42 & \textbf{63.7}\\
\bottomrule
\end{tabular*}
\end{table*}

\Cref{tab:main} combines PushObj and PushT evaluations. PushObj distinguishes archived baseline tasks (Reference) from the paired ECP task list (Paired); both cover the same seven shapes. On the paired list, our approach reaches 72.3\% with frozen dynamics and 76.9\% with AdaJEPA. On PushT, the adaptive configuration reaches 86\% clean success and 63.7\% averaged over six visual shifts, the highest mean among the compared controllers. With the augmented frozen model, ECP raises clean success from 72\% to 82\% and the shift mean from 46.0\% to 55.7\%. These gains transfer the PushObj correction settings unchanged; \cref{tab:pusht-ecp} reports all checkpoint and intervention comparisons.

\subsection{Manipulation and navigation}
\begin{table}[t]
\centering
\caption{\textbf{ECP improves both manipulation and navigation.} Success (\%) with at most 30 MPC replans; each row fixes the checkpoint and update rule. PushObj uses $S{+}P$. Maze uses $S$: movement-triggered for Frozen and loss-gated for AdaJEPA.}
\label{tab:domains}
\small
\setlength{\tabcolsep}{3pt}
\begin{tabular}{lllrr}
\toprule
Task & Checkpoint & Updates & Native & + ECP\\
\midrule
PushObj & Released & Frozen &43.7&\textbf{59.7}\\
& Released & AdaJEPA &69.4&\textbf{73.1}\\
& Data aug. & Frozen &61.1&\textbf{72.3}\\
& Data aug. & AdaJEPA &74.9&\textbf{76.9}\\
\midrule
Maze & Released & Frozen &48.0&\textbf{50.0}\\
& Released & AdaJEPA &63.0&\textbf{66.0}\\
\bottomrule
\end{tabular}
\end{table}

\paragraph{Improving both frozen and adaptive control.}
With the augmented-data checkpoint, ECP reaches 72.3\% frozen success and 76.9\% with AdaJEPA on PushObj (\cref{tab:domains}). Relative to the released-model baselines, the combined training and planning recipe improves success by 28.6 and 7.4 points, respectively. ECP also applies directly to the released checkpoint, raising frozen success from 43.7\% to 59.7\% and adaptive success from 69.4\% to 73.1\%. These paired comparisons demonstrate that execution feedback improves action selection under both fixed and adapting dynamics.

\paragraph{Recovering navigation with stagnation memory.}
On Maze, movement-triggered memory raises frozen success from 48\% to 50\%, recovering two tasks without changing a native success. With continued AdaJEPA updates, gated memory raises success from 63\% to 66\% and retains all 63 native successes. These improvements use the released checkpoint without new offline navigation training. \Cref{tab:maze} isolates the effect of correction type and activation schedule.
